\documentclass[11pt]{article}
\usepackage[margin=25mm]{geometry}
\usepackage{fontspec}
\setmainfont{Latin Modern Roman}
\usepackage{amsmath,amssymb,mathtools}
\usepackage{unicode-math}
\setmathfont{Latin Modern Math}
\usepackage{xurl,hyperref}
\hypersetup{colorlinks=true,urlcolor=blue}
\setcounter{secnumdepth}{0}
\setlength{\parindent}{0pt}
\setlength{\parskip}{5pt}
\setlength{\emergencystretch}{3em}
\newcommand{\sref}[2]{\hyperlink{source:#1:#2}{[S#2]}}
\newcommand{\cataloguescope}[1]{\paragraph{Recorded target.}#1}
\begin{document}
\section{Quasiperfect Numbers}\label{statement}

\subsection{Definitions and mathematical statement}

\paragraph{Shared notation.}$\mathbb N=\{1,2,\ldots\}$, and $\mathbb Z,\mathbb Q,\mathbb R,\mathbb C$ denote the integers, rationals, reals and complex numbers. Any other conventions are specified locally.


For a positive integer \(n\), write
\(\sigma(n)=\sum_{d\mid n,\,d>0}d\) for the sum of all its positive
divisors, including \(n\). A quasiperfect number is a positive integer
\(n\) satisfying \(\sigma(n)=2n+1\); equivalently its proper positive
divisors sum to \(n+1\).
\paragraph{Statement recorded in the supplied source.}
Does there exist \(n\in\mathbb N\) such that
\[
                             \sigma(n)=2n+1?
\] \sref{statement}{2}
\subsection{Short English statement}
Is any positive integer's divisor sum exactly one larger than twice that integer?
\subsection{Sources}
\begin{description}
\item[\hypertarget{source:statement:1}{S1}] ULAM.AI and the UnsolvedMath Contributors, \emph{UnsolvedMath: A Curated Collection of Open Mathematics Problems}, catalogue entry. CC BY 4.0. \url{https://huggingface.co/datasets/ulamai/UnsolvedMath}.
\item[\hypertarget{source:statement:2}{S2}] Akira Toyohara, Ye Tao and Siqiong Yao, \emph{Every quasiperfect number has at least eight distinct prime factors}, arXiv:2608.02066 (2026), Section 1.1, defining equation and existence question. Only the definition and question are used here. \url{https://arxiv.org/abs/2608.02066}.
\end{description}
\label{end:statement}
\end{document}
